% GENERATED FILE. Edit canonical lessons, then run npm run generate:books.
% canonical-source-hash: fnv1a64:dcf596ee1ea2d1b3
% canonical-lessons: UR-C127-unt, UR-C127-lomri, UR-C127-machchhar, UR-C127-makkhi, UR-C127-khargosh

\chapter{Camel, Fox, Mosquito, Fly, Rabbit}
\label{ch:ur-camel-fox-mosquito-fly-rabbit}
\hlchaptermodality{127}

\begin{chapteropening}
\textbf{By the end of this chapter:} \emph{I can say a camel, a fox, a mosquito, a fly and a rabbit.}

Complete the last lesson of chapter 127: I can say a camel, a fox, a mosquito, a fly and a rabbit.
\end{chapteropening}

\section[\texorpdfstring{\ur{اونٹ} (ūṅṭ)}{اونٹ (ūṅṭ)}]{\ur{اونٹ} (ūṅṭ) — a camel}

\label{lesson:UR-C127-unt}



\begin{quote}
Before the new one: say the Urdu for an ant, then the Urdu for a donkey.
\end{quote}

\subsection*{You'll want to know: \ur{اونٹ}}
\textbf{\ur{اونٹ}} — \emph{ūṅṭ} — \textquotedblleft{}a camel\textquotedblright{}.

\begin{cousinweb}
\textbf{\ur{اونٹ}} (\emph{ūṅṭ}) is a tadbhava, a word inherited from Sanskrit and reshaped by sound change. It comes from \emph{uṣṭra}, \textquotedblleft{}buffalo; camel\textquotedblright{} (CDIAL 2387).
\end{cousinweb}

\begin{grammarlens}[title={What you've built}]
One more word for this chapter.
\end{grammarlens}

\subsection*{Your turn}
\begin{itemize}
\raggedright
  \item \emph{Say it:} \emph{ūṅṭ}
  \item \emph{Say it:} \emph{ūṅṭ}, once more
  \item \emph{Say it:} say \emph{ūṅṭ}
  \item \emph{Recall:} say the Urdu for an ant, then the Urdu for a donkey, then say \emph{ūṅṭ} again
\end{itemize}

\begin{tcolorbox}[breakable,colback=teal!4,colframe=teal!35!black,title={Before you move on}]
What does \ur{اونٹ} mean? (\textquotedblleft{}A camel\textquotedblright{}.) Say it once more.
\end{tcolorbox}
\section[\texorpdfstring{\ur{لومڑی} (lomṛī)}{لومڑی (lomṛī)}]{\ur{لومڑی} (lomṛī) — a fox}

\label{lesson:UR-C127-lomri}



\begin{quote}
Before the new one: say the Urdu for a donkey, then the Urdu for a camel.
\end{quote}

\subsection*{You'll want to know: \ur{لومڑی}}
\textbf{\ur{لومڑی}} — \emph{lomṛī} — \textquotedblleft{}a fox\textquotedblright{}.

\begin{cousinweb}
\textbf{\ur{لومڑی}} (\emph{lomṛī}) is a tadbhava, a word inherited and reshaped by sound change. It comes from \emph{*lompaṭa}, \textquotedblleft{}fox\textquotedblright{}, a form CDIAL reconstructs from its descendants (CDIAL 11153). It is a variant CDIAL reconstructs beside Sanskrit \emph{lomaṭaka}, \textquotedblleft{}fox\textquotedblright{}; CDIAL suggests these forms may blend a word for \textquotedblleft{}fox\textquotedblright{} with words for \textquotedblleft{}hair\textquotedblright{} and \textquotedblleft{}tail\textquotedblright{}.
\end{cousinweb}

\begin{grammarlens}[title={What you've built}]
One more word for this chapter.
\end{grammarlens}

\subsection*{Your turn}
\begin{itemize}
\raggedright
  \item \emph{Say it:} \emph{lomṛī}
  \item \emph{Say it:} \emph{lomṛī}, once more
  \item \emph{Say it:} say \emph{lomṛī}
  \item \emph{Recall:} say the Urdu for a donkey, then the Urdu for a camel, then say \emph{lomṛī} again
\end{itemize}

\begin{tcolorbox}[breakable,colback=teal!4,colframe=teal!35!black,title={Before you move on}]
What does \ur{لومڑی} mean? (\textquotedblleft{}A fox\textquotedblright{}.) Say it once more.
\end{tcolorbox}
\section[\texorpdfstring{\ur{مچھر} (machchhar)}{مچھر (machchhar)}]{\ur{مچھر} (machchhar) — a mosquito}

\label{lesson:UR-C127-machchhar}



\begin{quote}
Before the new one: say the Urdu for a camel, then the Urdu for a fox.
\end{quote}

\subsection*{You'll want to know: \ur{مچھر}}
\textbf{\ur{مچھر}} — \emph{machchhar} — \textquotedblleft{}a mosquito\textquotedblright{}.

\begin{cousinweb}
\textbf{\ur{مچھر}} (\emph{machchhar}) is a tadbhava, a word inherited from Sanskrit and reshaped by sound change. It comes from \emph{matsara}, \textquotedblleft{}mosquito, fly\textquotedblright{} (CDIAL 9757). The Sanskrit word is known only from the old word-lists.
\end{cousinweb}

\begin{grammarlens}[title={What you've built}]
One more word for this chapter.
\end{grammarlens}

\subsection*{Your turn}
\begin{itemize}
\raggedright
  \item \emph{Say it:} \emph{machchhar}
  \item \emph{Say it:} \emph{machchhar}, once more
  \item \emph{Say it:} say \emph{machchhar}
  \item \emph{Recall:} say the Urdu for a camel, then the Urdu for a fox, then say \emph{machchhar} again
\end{itemize}

\begin{tcolorbox}[breakable,colback=teal!4,colframe=teal!35!black,title={Before you move on}]
What does \ur{مچھر} mean? (\textquotedblleft{}A mosquito\textquotedblright{}.) Say it once more.
\end{tcolorbox}
\section[\texorpdfstring{\ur{مکھی} (makkhī)}{مکھی (makkhī)}]{\ur{مکھی} (makkhī) — a fly}

\label{lesson:UR-C127-makkhi}



\begin{quote}
Before the new one: say the Urdu for a fox, then the Urdu for a mosquito.
\end{quote}

\subsection*{You'll want to know: \ur{مکھی}}
\textbf{\ur{مکھی}} — \emph{makkhī} — \textquotedblleft{}a fly\textquotedblright{}.

\begin{cousinweb}
\textbf{\ur{مکھی}} (\emph{makkhī}) is a tadbhava, a word inherited from Sanskrit and reshaped by sound change. It comes from \emph{makṣikā}, \textquotedblleft{}fly, bee\textquotedblright{} (CDIAL 9696).
\end{cousinweb}

\begin{grammarlens}[title={What you've built}]
One more word for this chapter.
\end{grammarlens}

\subsection*{Your turn}
\begin{itemize}
\raggedright
  \item \emph{Say it:} \emph{makkhī}
  \item \emph{Say it:} \emph{makkhī}, once more
  \item \emph{Say it:} say \emph{makkhī}
  \item \emph{Recall:} say the Urdu for a fox, then the Urdu for a mosquito, then say \emph{makkhī} again
\end{itemize}

\begin{tcolorbox}[breakable,colback=teal!4,colframe=teal!35!black,title={Before you move on}]
What does \ur{مکھی} mean? (\textquotedblleft{}A fly\textquotedblright{}.) Say it once more.
\end{tcolorbox}
\section[\texorpdfstring{\ur{خرگوش} (khargosh)}{خرگوش (khargosh)}]{\ur{خرگوش} (khargosh) — a rabbit}

\label{lesson:UR-C127-khargosh}



\begin{quote}
Before the new one: say the Urdu for a mosquito, then the Urdu for a fly.
\end{quote}

\subsection*{You'll want to know: \ur{خرگوش}}
\textbf{\ur{خرگوش}} — \emph{khargosh} — \textquotedblleft{}a rabbit\textquotedblright{}.

\begin{cousinweb}
\textbf{\ur{خرگوش}} (\emph{khargosh}) is a loan, usually traced to Persian \emph{khargōsh}, \textquotedblleft{}hare\textquotedblright{}. It is built from \emph{khar}, \textquotedblleft{}donkey\textquotedblright{}, and \emph{gōsh}, \textquotedblleft{}ear\textquotedblright{}.
\end{cousinweb}

\begin{grammarlens}[title={What you've built}]
One more word for this chapter.
\end{grammarlens}

\subsection*{Your turn}
\begin{itemize}
\raggedright
  \item \emph{Say it:} \emph{khargosh}
  \item \emph{Say it:} \emph{khargosh}, once more
  \item \emph{Say it:} say \emph{khargosh}
  \item \emph{Recall:} say the Urdu for a mosquito, then the Urdu for a fly, then say \emph{khargosh} again
\end{itemize}

\begin{tcolorbox}[breakable,colback=teal!4,colframe=teal!35!black,title={Before you move on}]
What does \ur{خرگوش} mean? (\textquotedblleft{}A rabbit\textquotedblright{}.) Say it once more.
\end{tcolorbox}
